\documentclass[11pt,a4paper]{article}
\usepackage{layout}
\begin{document}

\begin{canvas}
\pdfbookmark[0]{第一章 从开放协作到公共价值}{chapter1}
\node[anchor=north west,inner sep=0pt,outer sep=0pt] at (current page.north west){\includegraphics[width=210mm,height=297mm]{.build/assets/chapter1-cover.pdf}};
\Box{27}{227}{161}{\sffamily\bfseries\fontsize{11.5}{16}\selectfont\color{teal}AI 重塑开源，开源塑造 AI}
\end{canvas}

\begin{canvas}
\Header{1.1 · 双向重塑}{blue}\Footer{02}
\SectionTitle{blue}{1.1　人工智能时代的开源协作：双向重塑}{AI 改变创造与协作的节奏}{section11}
\Body{22}{66}{166}{AI 的影响不仅是增加代码产出，也开始改变工作如何发起、审阅和交付。智能体是能拆解任务并调用工具执行的 AI 系统；GitHub 的 Copilot 是 AI 编程助手。PR（合并请求）则是向项目提交修改、请求审阅与接纳的协作提案。}
\Box{22}{101}{80}{\sffamily\bfseries\fontsize{10.2}{14.5}\selectfont\color{blue}修改从哪里开始？}
\Box{108}{101}{80}{\sffamily\bfseries\fontsize{10.2}{14.5}\selectfont\color{teal}开发者如何使用助手？}
\Box{22}{109}{80}{\sffamily\fontsize{9}{13}\selectfont\color{navy}GitHub 中由智能体起草的 PR 占比}
\Box{108}{109}{80}{\sffamily\fontsize{9}{13}\selectfont\color{navy}Copilot 月活用户中使用智能体功能的比例}
\Box{22}{121}{80}{\sffamily\fontsize{8.4}{12}\selectfont\color{muted}2026 年 3 月}
\Rect{22}{128}{51}{3.5}{blue!8}
\Rect{22}{128}{8.67}{3.5}{blue}
\Box{76}{126}{27}{\sffamily\bfseries\fontsize{10}{14}\selectfont\color{blue}17\%}
\Box{22}{139}{80}{\sffamily\fontsize{8.4}{12}\selectfont\color{muted}2026 年 9 月}
\Rect{22}{146}{51}{3.5}{blue!8}
\Rect{22}{146}{20.400000000000002}{3.5}{blue}
\Box{76}{144}{27}{\sffamily\bfseries\fontsize{10}{14}\selectfont\color{blue}超过 40\%}
\LinkLine{42.400000000000006}{147.7}{46.400000000000006}{147.7}{blue}
\Box{22}{159}{80}{\sffamily\fontsize{7.8}{11}\selectfont\color{muted}0—100\%；“超过”数值以 40\% 为下界。}
\Box{108}{121}{80}{\sffamily\fontsize{8.4}{12}\selectfont\color{muted}报告所述上年}
\Rect{108}{128}{51}{3.5}{teal!8}
\Rect{108}{128}{17.849999999999998}{3.5}{teal}
\Box{162}{126}{27}{\sffamily\bfseries\fontsize{10}{14}\selectfont\color{teal}35\%}
\Box{108}{139}{80}{\sffamily\fontsize{8.4}{12}\selectfont\color{muted}2026 年报告}
\Rect{108}{146}{51}{3.5}{teal!8}
\Rect{108}{146}{28.05}{3.5}{teal}
\Box{162}{144}{27}{\sffamily\bfseries\fontsize{10}{14}\selectfont\color{teal}55\%}
\Box{108}{159}{80}{\sffamily\fontsize{7.8}{11}\selectfont\color{muted}分母是 Copilot 月活用户，不是全部开发者。}
\Rule{22}{172}{188}{line}
\Box{22}{178}{166}{\sffamily\bfseries\fontsize{10.2}{14.5}\selectfont\color{navy}从提出任务到长期维护：协作责任怎样衔接}
\Box{23}{193}{6}{\sffamily\fontsize{9.5}{14}\selectfont\color{ink}\includegraphics[width=6mm]{.build/assets/D01.pdf}}
\Box{32}{193}{32}{\sffamily\bfseries\fontsize{9.6}{14}\selectfont\color{blue}人定义问题}
\Box{22}{206}{38}{\sffamily\fontsize{8.5}{12}\selectfont\color{muted}明确需求与约束}
\LinkLine{59}{198}{63}{198}{muted!40}
\Box{66}{193}{6}{\sffamily\fontsize{9.5}{14}\selectfont\color{ink}\includegraphics[width=6mm]{.build/assets/D03.pdf}}
\Box{75}{193}{32}{\sffamily\bfseries\fontsize{9.6}{14}\selectfont\color{teal}智能体起草}
\Box{65}{206}{38}{\sffamily\fontsize{8.5}{12}\selectfont\color{muted}生成代码与修改}
\LinkLine{102}{198}{106}{198}{muted!40}
\Box{109}{193}{6}{\sffamily\fontsize{9.5}{14}\selectfont\color{ink}\includegraphics[width=6mm]{.build/assets/D09.pdf}}
\Box{118}{193}{32}{\sffamily\bfseries\fontsize{9.6}{14}\selectfont\color{gold}人审阅与接纳}
\Box{108}{206}{38}{\sffamily\fontsize{8.5}{12}\selectfont\color{muted}核对质量与风险}
\LinkLine{145}{198}{149}{198}{muted!40}
\Box{152}{193}{6}{\sffamily\fontsize{9.5}{14}\selectfont\color{ink}\includegraphics[width=6mm]{.build/assets/D10.pdf}}
\Box{161}{193}{32}{\sffamily\bfseries\fontsize{9.6}{14}\selectfont\color{violet}社区持续维护}
\Box{151}{206}{38}{\sffamily\fontsize{8.5}{12}\selectfont\color{muted}修复、更新与接续}
\Box{22}{218}{166}{\sffamily\fontsize{8.2}{12}\selectfont\color{muted}流程是解释示意：智能体可参与多个环节，人的审阅与项目治理仍需明确责任。}

\Body{22}{230}{166}{\textbf{需要观察的，是工作如何被组织。} 起草环节可以交给智能体，但需求是否合理、修改是否可靠、成果能否持续维护，仍需要项目规则与参与者共同判断。增长的活动量因此要与审阅、贡献和维护关系一起解读。}
\Note{22}{258}{166}{图 1-1　智能体进入协作流程。数据来源：\href{\GH}{GitHub Octoverse 2026}“Developers are reshaping their workflow”，2026-11-09 核对。GitHub 平台口径；两组比例的总体、窗口不同，不直接比较。2026 年 9 月属期后观察；流程为解释示意，数据不证明质量提升或 AI 的净因果效应。}
\end{canvas}

\begin{canvas}
\Header{1.1 · 中国模型与全球社区}{blue}\Footer{03}
\ContinueTitle{blue}{1.1 / 开源塑造 AI}{中国模型如何被全球继续使用}
\Subhead{59}{blue}{把“继续开发”“下载使用”与“在线调用”分开看}
\Body{22}{72}{166}{模型的影响不止来自发布者自身。其他开发者可以基于已有模型继续训练、调整用途或转换运行格式；使用者也可以下载模型，或通过服务平台调用它。下面三条证据分别回答不同问题。}
\shade[left color=blue!10,right color=blue!2,rounded corners=2mm,draw=blue!17,line width=.45pt] \Pos{22}{99} rectangle \Pos{188}{146};
\path[fill=blue!9] \Pos{81.75999999999999}{146} .. controls \Pos{134.88}{145} and \Pos{148.16}{141} .. \Pos{187}{142} -- \Pos{187}{145} -- \Pos{81.75999999999999}{146} -- cycle;
\Box{27}{102}{156}{\sffamily\bfseries\fontsize{10.2}{14.5}\selectfont\color{blue}Qwen｜社区在原有模型上继续开发}
\Box{27}{111}{57}{\sffamily\bfseries\fontsize{22}{27}\selectfont\color{blue}151,448}
\Box{27}{127}{57}{\sffamily\fontsize{8.7}{12}\selectfont\color{navy}HF 记录的衍生仓库}
\Box{91}{112}{91}{\sffamily\fontsize{9.3}{14.5}\selectfont\color{ink}其他开发者基于 Qwen 继续训练、调整用途或转换运行格式，形成可再次使用的版本。}
\Box{27}{136}{156}{\sffamily\fontsize{8}{11.5}\selectfont\color{muted}这些仓库体现再加工，不等于 15.1 万个从零训练的原创模型。}
\shade[left color=teal!10,right color=teal!2,rounded corners=2mm,draw=teal!17,line width=.45pt] \Pos{22}{152} rectangle \Pos{188}{199};
\path[fill=teal!9] \Pos{81.75999999999999}{199} .. controls \Pos{134.88}{198} and \Pos{148.16}{194} .. \Pos{187}{195} -- \Pos{187}{198} -- \Pos{81.75999999999999}{199} -- cycle;
\Box{27}{155}{156}{\sffamily\bfseries\fontsize{10.2}{14.5}\selectfont\color{teal}DeepSeek｜模型被下载到使用环境}
\Box{27}{164}{57}{\sffamily\bfseries\fontsize{22}{27}\selectfont\color{teal}47\%}
\Box{27}{180}{57}{\sffamily\fontsize{8.7}{12}\selectfont\color{navy}所追踪超大型模型下载份额}
\Box{91}{165}{91}{\sffamily\fontsize{9.3}{14.5}\selectfont\color{ink}ATOM 研究追踪的超大型开放语言模型累计下载中，DeepSeek 占近一半。下载可用于部署、测试或自动化流程。}
\Box{27}{189}{156}{\sffamily\fontsize{8}{11.5}\selectfont\color{muted}超大型按总参数≥2500 亿分组；参数是模型训练得到的数值，用于划分规模，不等于性能。}
\shade[left color=gold!10,right color=gold!2,rounded corners=2mm,draw=gold!17,line width=.45pt] \Pos{22}{205} rectangle \Pos{188}{252};
\path[fill=gold!9] \Pos{81.75999999999999}{252} .. controls \Pos{134.88}{251} and \Pos{148.16}{247} .. \Pos{187}{248} -- \Pos{187}{251} -- \Pos{81.75999999999999}{252} -- cycle;
\Box{27}{208}{156}{\sffamily\bfseries\fontsize{10.2}{14.5}\selectfont\color{gold}DeepSeek｜通过服务平台在线调用}
\Box{27}{217}{57}{\sffamily\bfseries\fontsize{22}{27}\selectfont\color{gold}31.1\%}
\Box{27}{233}{57}{\sffamily\fontsize{8.7}{12}\selectfont\color{navy}所选样本处理文本量份额}
\Box{91}{218}{91}{\sffamily\fontsize{9.3}{14.5}\selectfont\color{ink}在提供多模型调用服务的 OpenRouter 中，DeepSeek 占所选前十开放模型处理文本量的 31.1\%。}
\Box{27}{242}{156}{\sffamily\fontsize{8}{11.5}\selectfont\color{muted}以词元计量；词元是文本处理的基本单位。这是调用量信号，不是用户数或收入份额。}

\Note{22}{258}{166}{图 1-2a　中国模型的三种生态观察。来源：\href{\HF}{HF Summer 2026}；\href{\ATOM}{ATOM 论文 v1（2026-04）}图 9、12。HF 为夏季报告快照；下载为论文所追踪样本的累计值；调用为论文标注的 2026 年 1 月观察点。三者不是同一排行榜，也不是全球市场份额。}
\end{canvas}

\begin{canvas}
\Header{1.1 · 跨境接续与再创造}{blue}\Footer{04}
\ContinueTitle{blue}{1.1 / 从模型到数据与工具}{开放成果怎样进入他人的创新过程}
\Subhead{59}{blue}{DeepSeek：从模型输出，到训练材料与基础工具}
\Body{22}{72}{166}{一种模型的输出可以成为改进另一种模型的训练材料。NVIDIA 在 2025 年 5 月公开的数据说明中，列出由 DeepSeek-R1 生成的 1,212,994 条样本，并说明相应数据集用于训练 Llama-Nemotron 系列模型。}
\Box{22}{99}{166}{\sffamily\fontsize{8.6}{12}\selectfont\color{muted}历史案例 · NVIDIA 数据说明，2025 年 5 月}
\shade[left color=blue!10,right color=blue!2,rounded corners=2mm,draw=blue!17,line width=.45pt] \Pos{22}{108} rectangle \Pos{72}{152};
\path[fill=blue!9] \Pos{40.0}{152} .. controls \Pos{56.0}{151} and \Pos{60.0}{147} .. \Pos{71}{148} -- \Pos{71}{151} -- \Pos{40.0}{152} -- cycle;
\Box{26}{112}{6}{\sffamily\fontsize{9.5}{14}\selectfont\color{ink}\includegraphics[width=6mm]{.build/assets/D04.pdf}}
\Box{35}{112}{34}{\sffamily\bfseries\fontsize{9.6}{14}\selectfont\color{blue}DeepSeek-R1}
\Box{26}{126}{42}{\sffamily\fontsize{9}{13}\selectfont\color{navy}生成回答与推理过程}
\LinkLine{73}{129}{79}{129}{muted!60}
\shade[left color=teal!10,right color=teal!2,rounded corners=2mm,draw=teal!17,line width=.45pt] \Pos{80}{108} rectangle \Pos{130}{152};
\path[fill=teal!9] \Pos{98.0}{152} .. controls \Pos{114.0}{151} and \Pos{118.0}{147} .. \Pos{129}{148} -- \Pos{129}{151} -- \Pos{98.0}{152} -- cycle;
\Box{84}{112}{6}{\sffamily\fontsize{9.5}{14}\selectfont\color{ink}\includegraphics[width=6mm]{.build/assets/D05.pdf}}
\Box{93}{112}{34}{\sffamily\bfseries\fontsize{9.6}{14}\selectfont\color{teal}训练材料}
\Box{84}{126}{42}{\sffamily\fontsize{9}{13}\selectfont\color{navy}NVIDIA 整理并开放}
\Box{84}{138}{42}{\sffamily\bfseries\fontsize{9}{12}\selectfont\color{teal}其中 R1 样本约 121.3 万条}
\LinkLine{131}{129}{137}{129}{muted!60}
\shade[left color=gold!10,right color=gold!2,rounded corners=2mm,draw=gold!17,line width=.45pt] \Pos{138}{108} rectangle \Pos{188}{152};
\path[fill=gold!9] \Pos{156.0}{152} .. controls \Pos{172.0}{151} and \Pos{176.0}{147} .. \Pos{187}{148} -- \Pos{187}{151} -- \Pos{156.0}{152} -- cycle;
\Box{142}{112}{6}{\sffamily\fontsize{9.5}{14}\selectfont\color{ink}\includegraphics[width=6mm]{.build/assets/D04.pdf}}
\Box{151}{112}{34}{\sffamily\bfseries\fontsize{9.6}{14}\selectfont\color{gold}Llama-Nemotron}
\Box{142}{126}{42}{\sffamily\fontsize{9}{13}\selectfont\color{navy}使用相应数据集训练}

\Body{22}{158}{166}{这条路径复用的是模型生成的回答与推理过程，经整理后成为训练材料。它说明中国模型可以进入海外团队的再创造过程；不表示后续模型直接继承 DeepSeek 的参数，也不能据此分摊性能贡献。}
\Label{22}{186}{166}{模型之外，开放工具也在支撑共同研发}
\shade[left color=blue!10,right color=blue!2,rounded corners=2mm,draw=blue!17,line width=.45pt] \Pos{22}{196} rectangle \Pos{74}{231};
\path[fill=blue!9] \Pos{40.72}{231} .. controls \Pos{57.36}{230} and \Pos{61.52}{226} .. \Pos{73}{227} -- \Pos{73}{230} -- \Pos{40.72}{231} -- cycle;
\Box{25}{199}{46}{\sffamily\bfseries\fontsize{10.2}{14.5}\selectfont\color{blue}计算与推理}
\Box{25}{210}{46}{\sffamily\bfseries\fontsize{9}{13}\selectfont\color{navy}FlashMLA · DeepGEMM}
\Box{25}{221}{46}{\sffamily\fontsize{8.4}{12}\selectfont\color{muted}优化模型运算与推理执行}
\shade[left color=teal!10,right color=teal!2,rounded corners=2mm,draw=teal!17,line width=.45pt] \Pos{79}{196} rectangle \Pos{131}{231};
\path[fill=teal!9] \Pos{97.72}{231} .. controls \Pos{114.36}{230} and \Pos{118.52000000000001}{226} .. \Pos{130}{227} -- \Pos{130}{230} -- \Pos{97.72}{231} -- cycle;
\Box{82}{199}{46}{\sffamily\bfseries\fontsize{10.2}{14.5}\selectfont\color{teal}多芯片协作}
\Box{82}{210}{46}{\sffamily\bfseries\fontsize{9}{13}\selectfont\color{navy}DeepEP · DualPipe}
\Box{82}{221}{46}{\sffamily\fontsize{8.4}{12}\selectfont\color{muted}支持通信与并行训练调度}
\shade[left color=gold!10,right color=gold!2,rounded corners=2mm,draw=gold!17,line width=.45pt] \Pos{136}{196} rectangle \Pos{188}{231};
\path[fill=gold!9] \Pos{154.72}{231} .. controls \Pos{171.36}{230} and \Pos{175.52}{226} .. \Pos{187}{227} -- \Pos{187}{230} -- \Pos{154.72}{231} -- cycle;
\Box{139}{199}{46}{\sffamily\bfseries\fontsize{10.2}{14.5}\selectfont\color{gold}数据存储}
\Box{139}{210}{46}{\sffamily\bfseries\fontsize{9}{13}\selectfont\color{navy}3FS}
\Box{139}{221}{46}{\sffamily\fontsize{8.4}{12}\selectfont\color{muted}支撑训练数据的存取}

\Body{22}{240}{166}{模型、训练材料和基础工具是不同层次的贡献。把它们与实际采用者联系起来，才能进一步观察开放供给如何形成跨组织的体系能力。}
\Note{22}{260}{166}{图 1-2b　来源：\href{\NV}{NVIDIA 技术说明（2025-05-14）}表 2；\href{\DSINFRA}{DeepSeek 官方基础设施索引}（2025 年 2 月开源周）。两者均为主统计期之前的历史案例；工具为代表组件，不表示整个训练系统已完整开放。}
\end{canvas}

\begin{canvas}
\Header{1.2 · 发展目标}{teal}\Footer{05}
\SectionTitle{teal}{1.2　开源发展目标：从开放资源到公共能力}{明确共同方向、核对使用条件，再检验实际结果}{section12}
\Body{22}{66}{166}{联合国可持续发展目标（SDGs）提出 17 项共同发展目标。数字公共品（DPG）可包括开源软件、开放数据、人工智能系统与开放内容。数字公共品联盟（DPGA）通过九项要求评估方案的开放与设计条件；本地实施结果仍需独立检验。}
\Box{22}{101}{78}{\sffamily\bfseries\fontsize{10.6}{15}\selectfont\color{blue}17 项 SDGs｜共同发展方向}
\Box{108}{101}{80}{\sffamily\bfseries\fontsize{10.6}{15}\selectfont\color{teal}9 项 DPG 要求｜开放与设计}
\Rule{22}{111}{101}{blue!25}
\Rule{108}{111}{188}{teal!25}
\definecolor{sdg1}{HTML}{E5243B}
\Rect{22}{115.7}{0.9}{5.7}{sdg1}
\Box{24.5}{115.0}{7}{\sffamily\bfseries\fontsize{7.8}{10}\selectfont\color{muted}01}
\Box{31}{115.0}{30}{\sffamily\fontsize{8.3}{10.8}\selectfont\color{navy}无贫穷}
\definecolor{sdg2}{HTML}{DDA63A}
\Rect{22}{124.5}{0.9}{5.7}{sdg2}
\Box{24.5}{123.8}{7}{\sffamily\bfseries\fontsize{7.8}{10}\selectfont\color{muted}02}
\Box{31}{123.8}{30}{\sffamily\fontsize{8.3}{10.8}\selectfont\color{navy}零饥饿}
\definecolor{sdg3}{HTML}{4C9F38}
\Rect{22}{133.29999999999998}{0.9}{5.7}{sdg3}
\Box{24.5}{132.6}{7}{\sffamily\bfseries\fontsize{7.8}{10}\selectfont\color{muted}03}
\Box{31}{132.6}{30}{\sffamily\fontsize{8.3}{10.8}\selectfont\color{navy}良好健康与福祉}
\definecolor{sdg4}{HTML}{C5192D}
\Rect{22}{142.1}{0.9}{5.7}{sdg4}
\Box{24.5}{141.4}{7}{\sffamily\bfseries\fontsize{7.8}{10}\selectfont\color{muted}04}
\Box{31}{141.4}{30}{\sffamily\fontsize{8.3}{10.8}\selectfont\color{navy}优质教育}
\definecolor{sdg5}{HTML}{FF3A21}
\Rect{22}{150.89999999999998}{0.9}{5.7}{sdg5}
\Box{24.5}{150.2}{7}{\sffamily\bfseries\fontsize{7.8}{10}\selectfont\color{muted}05}
\Box{31}{150.2}{30}{\sffamily\fontsize{8.3}{10.8}\selectfont\color{navy}性别平等}
\definecolor{sdg6}{HTML}{26BDE2}
\Rect{22}{159.7}{0.9}{5.7}{sdg6}
\Box{24.5}{159.0}{7}{\sffamily\bfseries\fontsize{7.8}{10}\selectfont\color{muted}06}
\Box{31}{159.0}{30}{\sffamily\fontsize{8.3}{10.8}\selectfont\color{navy}清洁饮水和卫生设施}
\definecolor{sdg7}{HTML}{FCC30B}
\Rect{22}{168.5}{0.9}{5.7}{sdg7}
\Box{24.5}{167.8}{7}{\sffamily\bfseries\fontsize{7.8}{10}\selectfont\color{muted}07}
\Box{31}{167.8}{30}{\sffamily\fontsize{8.3}{10.8}\selectfont\color{navy}经济适用的清洁能源}
\definecolor{sdg8}{HTML}{A21942}
\Rect{22}{177.3}{0.9}{5.7}{sdg8}
\Box{24.5}{176.60000000000002}{7}{\sffamily\bfseries\fontsize{7.8}{10}\selectfont\color{muted}08}
\Box{31}{176.60000000000002}{30}{\sffamily\fontsize{8.3}{10.8}\selectfont\color{navy}体面工作和经济增长}
\definecolor{sdg9}{HTML}{FD6925}
\Rect{22}{186.1}{0.9}{5.7}{sdg9}
\Box{24.5}{185.4}{7}{\sffamily\bfseries\fontsize{7.8}{10}\selectfont\color{muted}09}
\Box{31}{185.4}{30}{\sffamily\fontsize{8.3}{10.8}\selectfont\color{navy}产业、创新和基础设施}
\definecolor{sdg10}{HTML}{DD1367}
\Rect{63}{115.7}{0.9}{5.7}{sdg10}
\Box{65.5}{115.0}{7}{\sffamily\bfseries\fontsize{7.8}{10}\selectfont\color{muted}10}
\Box{72}{115.0}{30}{\sffamily\fontsize{8.3}{10.8}\selectfont\color{navy}减少不平等}
\definecolor{sdg11}{HTML}{FD9D24}
\Rect{63}{124.5}{0.9}{5.7}{sdg11}
\Box{65.5}{123.8}{7}{\sffamily\bfseries\fontsize{7.8}{10}\selectfont\color{muted}11}
\Box{72}{123.8}{30}{\sffamily\fontsize{8.3}{10.8}\selectfont\color{navy}可持续城市和社区}
\definecolor{sdg12}{HTML}{BF8B2E}
\Rect{63}{133.29999999999998}{0.9}{5.7}{sdg12}
\Box{65.5}{132.6}{7}{\sffamily\bfseries\fontsize{7.8}{10}\selectfont\color{muted}12}
\Box{72}{132.6}{30}{\sffamily\fontsize{8.3}{10.8}\selectfont\color{navy}负责任消费和生产}
\definecolor{sdg13}{HTML}{3F7E44}
\Rect{63}{142.1}{0.9}{5.7}{sdg13}
\Box{65.5}{141.4}{7}{\sffamily\bfseries\fontsize{7.8}{10}\selectfont\color{muted}13}
\Box{72}{141.4}{30}{\sffamily\fontsize{8.3}{10.8}\selectfont\color{navy}气候行动}
\definecolor{sdg14}{HTML}{0A97D9}
\Rect{63}{150.89999999999998}{0.9}{5.7}{sdg14}
\Box{65.5}{150.2}{7}{\sffamily\bfseries\fontsize{7.8}{10}\selectfont\color{muted}14}
\Box{72}{150.2}{30}{\sffamily\fontsize{8.3}{10.8}\selectfont\color{navy}水下生物}
\definecolor{sdg15}{HTML}{56C02B}
\Rect{63}{159.7}{0.9}{5.7}{sdg15}
\Box{65.5}{159.0}{7}{\sffamily\bfseries\fontsize{7.8}{10}\selectfont\color{muted}15}
\Box{72}{159.0}{30}{\sffamily\fontsize{8.3}{10.8}\selectfont\color{navy}陆地生物}
\definecolor{sdg16}{HTML}{00689D}
\Rect{63}{168.5}{0.9}{5.7}{sdg16}
\Box{65.5}{167.8}{7}{\sffamily\bfseries\fontsize{7.8}{10}\selectfont\color{muted}16}
\Box{72}{167.8}{30}{\sffamily\fontsize{8.3}{10.8}\selectfont\color{navy}和平、正义与强大机构}
\definecolor{sdg17}{HTML}{19486A}
\Rect{63}{177.3}{0.9}{5.7}{sdg17}
\Box{65.5}{176.60000000000002}{7}{\sffamily\bfseries\fontsize{7.8}{10}\selectfont\color{muted}17}
\Box{72}{176.60000000000002}{30}{\sffamily\fontsize{8.3}{10.8}\selectfont\color{navy}促进目标实现的伙伴关系}
\Box{108}{115.0}{7}{\sffamily\bfseries\fontsize{8.1}{11}\selectfont\color{teal}01}
\Box{117}{115.0}{32}{\sffamily\bfseries\fontsize{8.6}{11}\selectfont\color{navy}SDG 相关性}
\Box{151}{115.0}{37}{\sffamily\fontsize{7.8}{10.7}\selectfont\color{muted}说明发展目标关联}
\Rule{108}{122.3}{188}{teal!12}
\Box{108}{123.8}{7}{\sffamily\bfseries\fontsize{8.1}{11}\selectfont\color{teal}02}
\Box{117}{123.8}{32}{\sffamily\bfseries\fontsize{8.6}{11}\selectfont\color{navy}开放许可}
\Box{151}{123.8}{37}{\sffamily\fontsize{7.8}{10.7}\selectfont\color{muted}允许合法复用}
\Rule{108}{131.1}{188}{teal!12}
\Box{108}{132.6}{7}{\sffamily\bfseries\fontsize{8.1}{11}\selectfont\color{teal}03}
\Box{117}{132.6}{32}{\sffamily\bfseries\fontsize{8.6}{11}\selectfont\color{navy}清晰权属}
\Box{151}{132.6}{37}{\sffamily\fontsize{7.8}{10.7}\selectfont\color{muted}说明权利归属}
\Rule{108}{139.9}{188}{teal!12}
\Box{108}{141.4}{7}{\sffamily\bfseries\fontsize{8.1}{11}\selectfont\color{teal}04}
\Box{117}{141.4}{32}{\sffamily\bfseries\fontsize{8.6}{11}\selectfont\color{navy}平台独立性}
\Box{151}{141.4}{37}{\sffamily\fontsize{7.8}{10.7}\selectfont\color{muted}减少封闭平台依赖}
\Rule{108}{148.70000000000002}{188}{teal!12}
\Box{108}{150.2}{7}{\sffamily\bfseries\fontsize{8.1}{11}\selectfont\color{teal}05}
\Box{117}{150.2}{32}{\sffamily\bfseries\fontsize{8.6}{11}\selectfont\color{navy}文档}
\Box{151}{150.2}{37}{\sffamily\fontsize{7.8}{10.7}\selectfont\color{muted}支持理解与使用}
\Rule{108}{157.5}{188}{teal!12}
\Box{108}{159.0}{7}{\sffamily\bfseries\fontsize{8.1}{11}\selectfont\color{teal}06}
\Box{117}{159.0}{32}{\sffamily\bfseries\fontsize{8.6}{11}\selectfont\color{navy}数据提取机制}
\Box{151}{159.0}{37}{\sffamily\fontsize{7.8}{10.7}\selectfont\color{muted}可提取非身份识别数据}
\Rule{108}{166.3}{188}{teal!12}
\Box{108}{167.8}{7}{\sffamily\bfseries\fontsize{8.1}{11}\selectfont\color{teal}07}
\Box{117}{167.8}{32}{\sffamily\bfseries\fontsize{8.6}{11}\selectfont\color{navy}隐私与法律}
\Box{151}{167.8}{37}{\sffamily\fontsize{7.8}{10.7}\selectfont\color{muted}遵守隐私及适用法律}
\Rule{108}{175.10000000000002}{188}{teal!12}
\Box{108}{176.60000000000002}{7}{\sffamily\bfseries\fontsize{8.1}{11}\selectfont\color{teal}08}
\Box{117}{176.60000000000002}{32}{\sffamily\bfseries\fontsize{8.6}{11}\selectfont\color{navy}标准与最佳实践}
\Box{151}{176.60000000000002}{37}{\sffamily\fontsize{7.8}{10.7}\selectfont\color{muted}遵循相关标准与实践}
\Rule{108}{183.90000000000003}{188}{teal!12}
\Box{108}{185.4}{7}{\sffamily\bfseries\fontsize{8.1}{11}\selectfont\color{teal}09}
\Box{117}{185.4}{32}{\sffamily\bfseries\fontsize{8.6}{11}\selectfont\color{navy}不造成伤害}
\Box{151}{185.4}{37}{\sffamily\fontsize{7.8}{10.7}\selectfont\color{muted}安全、内容与反骚扰保护}
\Box{22}{199}{166}{\sffamily\bfseries\fontsize{10.5}{15}\selectfont\color{navy}开放资源转化为公共能力，还要跨越四道实施门槛}
\draw[blue!35,line width=0.9pt,-{Stealth[length=1.4mm]}] \Pos{47}{219} .. controls \Pos{58}{216} and \Pos{63}{222} .. \Pos{76}{219};
\fill[blue!5] \Pos{40}{219} circle[radius=6.7mm];
\shade[top color=white,bottom color=blue!28,shading angle=30,draw=blue!35,line width=.45pt] \Pos{40}{219} circle[radius=5.5mm];
\draw[white,line width=.7pt] \Pos{36.81}{215.81} arc[start angle=135,end angle=65,radius=4.51mm];
\Box{37.14}{216.14}{5.720000000000001}{\sffamily\fontsize{9.5}{14}\selectfont\color{ink}\includegraphics[width=5.720000000000001mm]{.build/assets/D05.pdf}}
\Box{22}{227}{36}{\sffamily\bfseries\fontsize{9.8}{14}\selectfont\color{blue}\centering 可获得}
\Box{20}{236}{40}{\sffamily\fontsize{8}{11}\selectfont\color{muted}\centering 版本 · 许可 · 文档}
\draw[teal!35,line width=0.9pt,-{Stealth[length=1.4mm]}] \Pos{90}{219} .. controls \Pos{101}{216} and \Pos{106}{222} .. \Pos{119}{219};
\fill[teal!5] \Pos{83}{219} circle[radius=6.7mm];
\shade[top color=white,bottom color=teal!28,shading angle=30,draw=teal!35,line width=.45pt] \Pos{83}{219} circle[radius=5.5mm];
\draw[white,line width=.7pt] \Pos{79.81}{215.81} arc[start angle=135,end angle=65,radius=4.51mm];
\Box{80.14}{216.14}{5.720000000000001}{\sffamily\fontsize{9.5}{14}\selectfont\color{ink}\includegraphics[width=5.720000000000001mm]{.build/assets/D08.pdf}}
\Box{65}{227}{36}{\sffamily\bfseries\fontsize{9.8}{14}\selectfont\color{teal}\centering 可复用}
\Box{63}{236}{40}{\sffamily\fontsize{8}{11}\selectfont\color{muted}\centering 成本 · 语言 · 工具}
\draw[gold!35,line width=0.9pt,-{Stealth[length=1.4mm]}] \Pos{133}{219} .. controls \Pos{144}{216} and \Pos{149}{222} .. \Pos{162}{219};
\fill[gold!5] \Pos{126}{219} circle[radius=6.7mm];
\shade[top color=white,bottom color=gold!28,shading angle=30,draw=gold!35,line width=.45pt] \Pos{126}{219} circle[radius=5.5mm];
\draw[white,line width=.7pt] \Pos{122.81}{215.81} arc[start angle=135,end angle=65,radius=4.51mm];
\Box{123.14}{216.14}{5.720000000000001}{\sffamily\fontsize{9.5}{14}\selectfont\color{ink}\includegraphics[width=5.720000000000001mm]{.build/assets/D09.pdf}}
\Box{108}{227}{36}{\sffamily\bfseries\fontsize{9.8}{14}\selectfont\color{gold}\centering 可维护}
\Box{106}{236}{40}{\sffamily\fontsize{8}{11}\selectfont\color{muted}\centering 责任 · 资金 · 安全}
\fill[violet!5] \Pos{169}{219} circle[radius=6.7mm];
\shade[top color=white,bottom color=violet!28,shading angle=30,draw=violet!35,line width=.45pt] \Pos{169}{219} circle[radius=5.5mm];
\draw[white,line width=.7pt] \Pos{165.81}{215.81} arc[start angle=135,end angle=65,radius=4.51mm];
\Box{166.14}{216.14}{5.720000000000001}{\sffamily\fontsize{9.5}{14}\selectfont\color{ink}\includegraphics[width=5.720000000000001mm]{.build/assets/D10.pdf}}
\Box{151}{227}{36}{\sffamily\bfseries\fontsize{9.8}{14}\selectfont\color{violet}\centering 有成效}
\Box{149}{236}{40}{\sffamily\fontsize{8}{11}\selectfont\color{muted}\centering 覆盖 · 质量 · 公平}

\Body{22}{247}{166}{AI 普惠需要本地伙伴参与问题定义与技术改进，也需要基金会、企业、研究机构和社区共同承担长期维护。资源开放之后，谁能够使用、谁负责维护、谁持续受益，应成为公共投入的共同问题。}
\Note{22}{268}{166}{图 1-3　来源：\href{\UNGoals}{联合国 SDGs}、\href{\DPG}{DPGA 标准}（2026-11-09 核对；中文简述与改绘部分 CC BY-SA 4.0）。两套体系不逐项对应；四道门槛为本报告分析框架，非自动生效的因果链。}
\end{canvas}

\begin{canvas}
\Header{1.2 · 发展目标}{teal}\Footer{06}
\ContinueTitle{teal}{1.2 / 开源发展的共同语言}{OSDGs：四个支柱共同支撑公共价值}
\Subhead{59}{teal}{不仅看增长，也看谁参与、谁接续、谁受益}
\Body{22}{72}{166}{本报告提出开源发展目标（OSDGs），把全球开源的变化放回发展问题：人才是否持续成长，资源是否可以复用，成果是否形成连接，治理是否支持共同受益。四个支柱互为补充，十个方向细化观察重点。}
\draw[blue!12,line width=5pt] \Pos{79}{125} .. controls \Pos{88}{125} and \Pos{84}{135} .. \Pos{91}{142};
\draw[blue!55,line width=0.7pt] \Pos{79}{125} .. controls \Pos{88}{125} and \Pos{84}{135} .. \Pos{91}{142};
\draw[teal!12,line width=5pt] \Pos{131}{125} .. controls \Pos{122}{125} and \Pos{127}{135} .. \Pos{119}{142};
\draw[teal!55,line width=0.7pt] \Pos{131}{125} .. controls \Pos{122}{125} and \Pos{127}{135} .. \Pos{119}{142};
\draw[gold!12,line width=5pt] \Pos{79}{184} .. controls \Pos{88}{184} and \Pos{84}{172} .. \Pos{91}{165};
\draw[gold!55,line width=0.7pt] \Pos{79}{184} .. controls \Pos{88}{184} and \Pos{84}{172} .. \Pos{91}{165};
\draw[violet!12,line width=5pt] \Pos{131}{184} .. controls \Pos{122}{184} and \Pos{127}{172} .. \Pos{119}{165};
\draw[violet!55,line width=0.7pt] \Pos{131}{184} .. controls \Pos{122}{184} and \Pos{127}{172} .. \Pos{119}{165};
\Box{86}{135}{38}{\sffamily\fontsize{9.5}{14}\selectfont\color{ink}\includegraphics[width=38mm]{.build/assets/osdgs-display.pdf}}
\Box{79}{177}{52}{\sffamily\bfseries\fontsize{9.1}{13}\selectfont\color{teal}\centering 共同发展 · 全球接续}
\shade[left color=blue!10,right color=blue!2,rounded corners=2mm,draw=blue!17,line width=.45pt] \Pos{22}{99} rectangle \Pos{80}{152};
\path[fill=blue!9] \Pos{42.879999999999995}{152} .. controls \Pos{61.440000000000005}{151} and \Pos{66.08}{147} .. \Pos{79}{148} -- \Pos{79}{151} -- \Pos{42.879999999999995}{152} -- cycle;
\Box{25}{102}{52}{\sffamily\bfseries\fontsize{9.6}{14}\selectfont\color{blue}X1 人才规模与多样性}
\Box{25}{111}{52}{\sffamily\fontsize{8.4}{12}\selectfont\color{muted}扩大参与与持续成长}
\Box{25}{123.0}{4.2}{\sffamily\fontsize{9.5}{14}\selectfont\color{ink}\includegraphics[width=4.2mm]{.build/assets/D01.pdf}}
\Box{32}{123.0}{45}{\sffamily\fontsize{8.6}{12}\selectfont\color{ink}D01 规模与密度}
\Box{25}{131.7}{4.2}{\sffamily\fontsize{9.5}{14}\selectfont\color{ink}\includegraphics[width=4.2mm]{.build/assets/D02.pdf}}
\Box{32}{131.7}{45}{\sffamily\fontsize{8.6}{12}\selectfont\color{ink}D02 生态多样性}
\shade[left color=teal!10,right color=teal!2,rounded corners=2mm,draw=teal!17,line width=.45pt] \Pos{130}{99} rectangle \Pos{188}{152};
\path[fill=teal!9] \Pos{150.88}{152} .. controls \Pos{169.44}{151} and \Pos{174.07999999999998}{147} .. \Pos{187}{148} -- \Pos{187}{151} -- \Pos{150.88}{152} -- cycle;
\Box{133}{102}{52}{\sffamily\bfseries\fontsize{9.6}{14}\selectfont\color{teal}X2 技术产出与开放供给}
\Box{133}{111}{52}{\sffamily\fontsize{8.4}{12}\selectfont\color{muted}形成可复用的开放资源}
\Box{133}{123.0}{4.2}{\sffamily\fontsize{9.5}{14}\selectfont\color{ink}\includegraphics[width=4.2mm]{.build/assets/D03.pdf}}
\Box{140}{123.0}{45}{\sffamily\fontsize{8.6}{12}\selectfont\color{ink}D03 代码产出}
\Box{133}{131.7}{4.2}{\sffamily\fontsize{9.5}{14}\selectfont\color{ink}\includegraphics[width=4.2mm]{.build/assets/D04.pdf}}
\Box{140}{131.7}{45}{\sffamily\fontsize{8.6}{12}\selectfont\color{ink}D04 模型开放供给}
\Box{133}{140.4}{4.2}{\sffamily\fontsize{9.5}{14}\selectfont\color{ink}\includegraphics[width=4.2mm]{.build/assets/D05.pdf}}
\Box{140}{140.4}{45}{\sffamily\fontsize{8.6}{12}\selectfont\color{ink}D05 数据开放供给}
\shade[left color=gold!10,right color=gold!2,rounded corners=2mm,draw=gold!17,line width=.45pt] \Pos{22}{160} rectangle \Pos{80}{213};
\path[fill=gold!9] \Pos{42.879999999999995}{213} .. controls \Pos{61.440000000000005}{212} and \Pos{66.08}{208} .. \Pos{79}{209} -- \Pos{79}{212} -- \Pos{42.879999999999995}{213} -- cycle;
\Box{25}{163}{52}{\sffamily\bfseries\fontsize{9.6}{14}\selectfont\color{gold}X3 全球连接与网络效应}
\Box{25}{172}{52}{\sffamily\fontsize{8.4}{12}\selectfont\color{muted}让成果跨境接续生长}
\Box{25}{184.0}{4.2}{\sffamily\fontsize{9.5}{14}\selectfont\color{ink}\includegraphics[width=4.2mm]{.build/assets/D06.pdf}}
\Box{32}{184.0}{45}{\sffamily\fontsize{8.6}{12}\selectfont\color{ink}D06 代码影响力}
\Box{25}{192.7}{4.2}{\sffamily\fontsize{9.5}{14}\selectfont\color{ink}\includegraphics[width=4.2mm]{.build/assets/D07.pdf}}
\Box{32}{192.7}{45}{\sffamily\fontsize{8.6}{12}\selectfont\color{ink}D07 模型影响力}
\Box{25}{201.4}{4.2}{\sffamily\fontsize{9.5}{14}\selectfont\color{ink}\includegraphics[width=4.2mm]{.build/assets/D08.pdf}}
\Box{32}{201.4}{45}{\sffamily\fontsize{8.6}{12}\selectfont\color{ink}D08 数据集影响力}
\shade[left color=violet!10,right color=violet!2,rounded corners=2mm,draw=violet!17,line width=.45pt] \Pos{130}{160} rectangle \Pos{188}{213};
\path[fill=violet!9] \Pos{150.88}{213} .. controls \Pos{169.44}{212} and \Pos{174.07999999999998}{208} .. \Pos{187}{209} -- \Pos{187}{212} -- \Pos{150.88}{213} -- cycle;
\Box{133}{163}{52}{\sffamily\bfseries\fontsize{9.6}{14}\selectfont\color{violet}X4 开源善治与公共价值}
\Box{133}{172}{52}{\sffamily\fontsize{8.4}{12}\selectfont\color{muted}让协作持续、成果共享}
\Box{133}{184.0}{4.2}{\sffamily\fontsize{9.5}{14}\selectfont\color{ink}\includegraphics[width=4.2mm]{.build/assets/D09.pdf}}
\Box{140}{184.0}{45}{\sffamily\fontsize{8.6}{12}\selectfont\color{ink}D09 治理成熟度}
\Box{133}{192.7}{4.2}{\sffamily\fontsize{9.5}{14}\selectfont\color{ink}\includegraphics[width=4.2mm]{.build/assets/D10.pdf}}
\Box{140}{192.7}{45}{\sffamily\fontsize{8.6}{12}\selectfont\color{ink}D10 公共福祉}
\Box{22}{218}{166}{\sffamily\fontsize{8.2}{12}\selectfont\color{muted}\centering OSDGs / Open Source Development Goals}

\Cap{227}{图 1-4　开源发展目标：四大战略支柱与十个观察维度}
\Body{22}{237}{166}{\textbf{人才与供给相互支撑。} 参与机会需要转化为持续贡献，开放成果需要具备可理解、可使用的条件。}
\Body{22}{251}{166}{\textbf{连接与治理共同延长价值。} 把跨境接续、共同维护与受益群体纳入观察，才能解释资源增长之后的体系能力。}
\Note{22}{270}{166}{来源：本报告获批框架与动态星座 v2.0.0。OSDGs 是本报告研究框架，不是联合国目标或认证；十维编号用于导视，不是测算权重。}
\end{canvas}

\begin{canvas}
\Header{1.3 · 观察方法}{gold}\Footer{07}
\SectionTitle{gold}{1.3　观察方法：跨平台的全球开源全景}{OSDI：把发展目标转化为具体问题与观察信号}{section13}
\Body{22}{66}{166}{开源发展指标体系（OSDI）连接四个发展目标、十个观察方向与 29 项候选指标。目标说明“希望改善什么”，方向划分“从哪里观察”，指标回答“用什么证据判断变化”。下面以人才目标的一条分支为例。}
\Box{22}{101}{166}{\sffamily\bfseries\fontsize{10.2}{14.5}\selectfont\color{navy}示例：关注人才发展，为什么不能只看新增账号？}
\shade[left color=blue!10,right color=blue!2,rounded corners=2mm,draw=blue!17,line width=.45pt] \Pos{22}{114} rectangle \Pos{56}{164};
\path[fill=blue!9] \Pos{34.24}{164} .. controls \Pos{45.120000000000005}{163} and \Pos{47.84}{159} .. \Pos{55}{160} -- \Pos{55}{163} -- \Pos{34.24}{164} -- cycle;
\Box{26}{118}{26}{\sffamily\fontsize{8.5}{12}\selectfont\color{muted}一级 · 目标}
\Box{26}{128}{26}{\sffamily\bfseries\fontsize{10.4}{16}\selectfont\color{blue}X1 人才规模与多样性}
\Box{26}{150}{26}{\sffamily\fontsize{8.3}{12}\selectfont\color{ink}培养持续贡献者}
\shade[left color=teal!10,right color=teal!2,rounded corners=2mm,draw=teal!17,line width=.45pt] \Pos{62}{114} rectangle \Pos{96}{164};
\path[fill=teal!9] \Pos{74.24}{164} .. controls \Pos{85.12}{163} and \Pos{87.84}{159} .. \Pos{95}{160} -- \Pos{95}{163} -- \Pos{74.24}{164} -- cycle;
\Box{66}{118}{26}{\sffamily\fontsize{8.5}{12}\selectfont\color{muted}二级 · 方向}
\Box{66}{128}{26}{\sffamily\bfseries\fontsize{10.4}{16}\selectfont\color{teal}D01 规模与密度}
\Box{66}{150}{26}{\sffamily\fontsize{8.3}{12}\selectfont\color{ink}观察规模与分布}
\LinkLine{56}{139}{61}{139}{blue!50}
\LinkLine{96}{139}{102}{139}{teal!50}
\shade[left color=blue!10,right color=blue!2,rounded corners=2mm,draw=blue!17,line width=.45pt] \Pos{103}{113} rectangle \Pos{188}{129};
\path[fill=blue!9] \Pos{133.6}{129} .. controls \Pos{160.8}{128} and \Pos{167.6}{124} .. \Pos{187}{125} -- \Pos{187}{128} -- \Pos{133.6}{129} -- cycle;
\Box{106}{115}{60}{\sffamily\bfseries\fontsize{9.4}{13}\selectfont\color{navy}活跃开发者规模}
\Box{167}{115}{18}{\sffamily\bfseries\fontsize{8.6}{12}\selectfont\color{blue}北极星}
\Box{106}{122}{77}{\sffamily\fontsize{8.6}{12}\selectfont\color{muted}有多少人在实际参与协作？}
\shade[left color=teal!10,right color=teal!2,rounded corners=2mm,draw=teal!17,line width=.45pt] \Pos{103}{131} rectangle \Pos{188}{147};
\path[fill=teal!9] \Pos{133.6}{147} .. controls \Pos{160.8}{146} and \Pos{167.6}{142} .. \Pos{187}{143} -- \Pos{187}{146} -- \Pos{133.6}{147} -- cycle;
\Box{106}{133}{60}{\sffamily\bfseries\fontsize{9.4}{13}\selectfont\color{navy}人才聚集密度}
\Box{167}{133}{18}{\sffamily\bfseries\fontsize{8.6}{12}\selectfont\color{teal}诊断}
\Box{106}{140}{77}{\sffamily\fontsize{8.6}{12}\selectfont\color{muted}参与者集中在哪里？}
\shade[left color=gold!10,right color=gold!2,rounded corners=2mm,draw=gold!17,line width=.45pt] \Pos{103}{149} rectangle \Pos{188}{165};
\path[fill=gold!9] \Pos{133.6}{165} .. controls \Pos{160.8}{164} and \Pos{167.6}{160} .. \Pos{187}{161} -- \Pos{187}{164} -- \Pos{133.6}{165} -- cycle;
\Box{106}{151}{60}{\sffamily\bfseries\fontsize{9.4}{13}\selectfont\color{navy}新晋开发者增速}
\Box{167}{151}{18}{\sffamily\bfseries\fontsize{8.6}{12}\selectfont\color{gold}驱动}
\Box{106}{158}{77}{\sffamily\fontsize{8.6}{12}\selectfont\color{muted}新的参与力量怎样进入？}

\Note{22}{173}{166}{图 1-5　目标、方向与指标的关系示例。“北极星”提示长期方向，“诊断”解释差异，“驱动”跟踪变化；三类是用途，不是指标权重或证据等级。完整目录见下一页。}
\Label{22}{189}{166}{指标怎样服务“从洞察到共建”}
\Rect{22}{199}{166}{7}{paper}
\Box{25}{200}{30}{\sffamily\bfseries\fontsize{8.8}{12}\selectfont\color{navy}使用环节}
\Box{59}{200}{65}{\sffamily\bfseries\fontsize{8.8}{12}\selectfont\color{navy}指标帮助回答什么}
\Box{127}{200}{58}{\sffamily\bfseries\fontsize{8.8}{12}\selectfont\color{navy}行动前还需核对}
\Box{25}{209}{31}{\sffamily\bfseries\fontsize{10.2}{14.5}\selectfont\color{blue}透视与诊断}
\Box{59}{209}{64}{\sffamily\fontsize{9}{13}\selectfont\color{ink}哪里供给丰富，哪里连接薄弱？}
\Box{127}{209}{58}{\sffamily\fontsize{8.7}{12.5}\selectfont\color{muted}同窗口、同对象、可比基线}
\Rule{22}{219.5}{188}{line}
\Box{25}{222}{31}{\sffamily\bfseries\fontsize{10.2}{14.5}\selectfont\color{teal}引导与激励}
\Box{59}{222}{64}{\sffamily\fontsize{9}{13}\selectfont\color{ink}如何支持持续贡献与共同维护？}
\Box{127}{222}{58}{\sffamily\fontsize{8.7}{12.5}\selectfont\color{muted}贡献质量、角色与激励后果}
\Rule{22}{232.5}{188}{line}
\Box{25}{235}{31}{\sffamily\bfseries\fontsize{10.2}{14.5}\selectfont\color{gold}赋能与决策}
\Box{59}{235}{64}{\sffamily\fontsize{9}{13}\selectfont\color{ink}公共投入和合作应优先补哪里？}
\Box{127}{235}{58}{\sffamily\fontsize{8.7}{12.5}\selectfont\color{muted}实际成本、本地需求与承接能力}
\Rule{22}{245.5}{188}{line}
\Box{25}{248}{31}{\sffamily\bfseries\fontsize{10.2}{14.5}\selectfont\color{violet}连接与共建}
\Box{59}{248}{64}{\sffamily\fontsize{9}{13}\selectfont\color{ink}哪些伙伴可以共同建设公共能力？}
\Box{127}{248}{58}{\sffamily\fontsize{8.7}{12.5}\selectfont\color{muted}开放条件、责任安排与受益结果}
\Rule{22}{258.5}{188}{line}

\Note{22}{266}{166}{来源：获批第一章指标框架与本轮用户补充材料，按决策阅读重绘。指标定义、计算窗口与可得性仍须逐项核定；不能直接合成为国家综合实力或公共成效总分。}
\end{canvas}

\begin{canvas}
\Header{1.3 · OSDI 指标目录}{gold}\Footer{08}
\ContinueTitle{gold}{1.3 / 由目标到可观察的信号}{29 项候选指标，分别观察什么}
\Body{22}{62}{166}{下表完整保留四目标、十方向与 29 项候选指标。名称中的“高质量”“采用”“合规”等仍需操作性定义；目录不表示全部已经测量。“目标关联”沿用补充材料，仅指研究意图。}
\Rect{22}{86}{166}{8}{navy}
\Box{25}{87}{36}{\sffamily\bfseries\fontsize{8.8}{12}\selectfont\color{white}方向}
\Box{62}{87}{75}{\sffamily\bfseries\fontsize{8.8}{12}\selectfont\color{white}候选指标}
\Box{139}{87}{16}{\sffamily\bfseries\fontsize{8.8}{12}\selectfont\color{white}用途}
\Box{158}{87}{29}{\sffamily\bfseries\fontsize{8.8}{12}\selectfont\color{white}目标关联}
\Rect{22}{95}{166}{6.5}{blue!9}
\Box{25}{96}{160}{\sffamily\bfseries\fontsize{9.2}{12.5}\selectfont\color{blue}X1 人才规模与多样性}
\Box{25}{102.3}{34}{\sffamily\bfseries\fontsize{8.3}{11.3}\selectfont\color{navy}D01 规模与密度}
\Box{62}{102.3}{73}{\sffamily\fontsize{8.7}{11.5}\selectfont\color{ink}活跃开发者规模}
\Box{140}{102.3}{14}{\sffamily\bfseries\fontsize{8.6}{11.5}\selectfont\color{blue}北极星}
\Box{158}{102.3}{27}{\sffamily\fontsize{8.5}{11.5}\selectfont\color{muted}8、9}
\Rule{60}{106.4}{188}{line}
\Box{62}{106.95}{73}{\sffamily\fontsize{8.7}{11.5}\selectfont\color{ink}人才聚集密度}
\Box{140}{106.95}{14}{\sffamily\bfseries\fontsize{8.6}{11.5}\selectfont\color{blue}诊断}
\Box{158}{106.95}{27}{\sffamily\fontsize{8.5}{11.5}\selectfont\color{muted}9、11}
\Rule{60}{111.05000000000001}{188}{line}
\Box{62}{111.60000000000001}{73}{\sffamily\fontsize{8.7}{11.5}\selectfont\color{ink}新晋开发者增速}
\Box{140}{111.60000000000001}{14}{\sffamily\bfseries\fontsize{8.6}{11.5}\selectfont\color{blue}驱动}
\Box{158}{111.60000000000001}{27}{\sffamily\fontsize{8.5}{11.5}\selectfont\color{muted}4、8}
\Rule{60}{115.70000000000002}{188}{line}
\Box{25}{116.25000000000001}{34}{\sffamily\bfseries\fontsize{8.3}{11.3}\selectfont\color{navy}D02 生态多样性}
\Box{62}{116.25000000000001}{73}{\sffamily\fontsize{8.7}{11.5}\selectfont\color{ink}全球开源协作占比}
\Box{140}{116.25000000000001}{14}{\sffamily\bfseries\fontsize{8.6}{11.5}\selectfont\color{blue}诊断}
\Box{158}{116.25000000000001}{27}{\sffamily\fontsize{8.5}{11.5}\selectfont\color{muted}10、17}
\Rule{60}{120.35000000000002}{188}{line}
\Box{62}{120.90000000000002}{73}{\sffamily\fontsize{8.7}{11.5}\selectfont\color{ink}学术机构贡献占比}
\Box{140}{120.90000000000002}{14}{\sffamily\bfseries\fontsize{8.6}{11.5}\selectfont\color{blue}诊断}
\Box{158}{120.90000000000002}{27}{\sffamily\fontsize{8.5}{11.5}\selectfont\color{muted}4、9}
\Rule{60}{125.00000000000003}{188}{line}
\Rect{22}{125.25000000000003}{166}{6.5}{teal!9}
\Box{25}{126.25000000000003}{160}{\sffamily\bfseries\fontsize{9.2}{12.5}\selectfont\color{teal}X2 技术产出与开放供给}
\Box{25}{132.55000000000004}{34}{\sffamily\bfseries\fontsize{8.3}{11.3}\selectfont\color{navy}D03 代码产出}
\Box{62}{132.55000000000004}{73}{\sffamily\fontsize{8.7}{11.5}\selectfont\color{ink}有效活跃仓库规模}
\Box{140}{132.55000000000004}{14}{\sffamily\bfseries\fontsize{8.6}{11.5}\selectfont\color{teal}北极星}
\Box{158}{132.55000000000004}{27}{\sffamily\fontsize{8.5}{11.5}\selectfont\color{muted}9}
\Rule{60}{136.65000000000003}{188}{line}
\Box{62}{137.20000000000005}{73}{\sffamily\fontsize{8.7}{11.5}\selectfont\color{ink}代码协作贡献度}
\Box{140}{137.20000000000005}{14}{\sffamily\bfseries\fontsize{8.6}{11.5}\selectfont\color{teal}北极星}
\Box{158}{137.20000000000005}{27}{\sffamily\fontsize{8.5}{11.5}\selectfont\color{muted}9、17}
\Rule{60}{141.30000000000004}{188}{line}
\Box{62}{141.85000000000005}{73}{\sffamily\fontsize{8.7}{11.5}\selectfont\color{ink}高质量 PR 合入量}
\Box{140}{141.85000000000005}{14}{\sffamily\bfseries\fontsize{8.6}{11.5}\selectfont\color{teal}驱动}
\Box{158}{141.85000000000005}{27}{\sffamily\fontsize{8.5}{11.5}\selectfont\color{muted}9}
\Rule{60}{145.95000000000005}{188}{line}
\Box{25}{146.50000000000006}{34}{\sffamily\bfseries\fontsize{8.3}{11.3}\selectfont\color{navy}D04 模型开放供给}
\Box{62}{146.50000000000006}{73}{\sffamily\fontsize{8.7}{11.5}\selectfont\color{ink}开源模型发布量}
\Box{140}{146.50000000000006}{14}{\sffamily\bfseries\fontsize{8.6}{11.5}\selectfont\color{teal}驱动}
\Box{158}{146.50000000000006}{27}{\sffamily\fontsize{8.5}{11.5}\selectfont\color{muted}9}
\Rule{60}{150.60000000000005}{188}{line}
\Box{62}{151.15000000000006}{73}{\sffamily\fontsize{8.7}{11.5}\selectfont\color{ink}模型供给方多样性}
\Box{140}{151.15000000000006}{14}{\sffamily\bfseries\fontsize{8.6}{11.5}\selectfont\color{teal}诊断}
\Box{158}{151.15000000000006}{27}{\sffamily\fontsize{8.5}{11.5}\selectfont\color{muted}9、10}
\Rule{60}{155.25000000000006}{188}{line}
\Box{25}{155.80000000000007}{34}{\sffamily\bfseries\fontsize{8.3}{11.3}\selectfont\color{navy}D05 数据开放供给}
\Box{62}{155.80000000000007}{73}{\sffamily\fontsize{8.7}{11.5}\selectfont\color{ink}开源数据集发布量}
\Box{140}{155.80000000000007}{14}{\sffamily\bfseries\fontsize{8.6}{11.5}\selectfont\color{teal}驱动}
\Box{158}{155.80000000000007}{27}{\sffamily\fontsize{8.5}{11.5}\selectfont\color{muted}9}
\Rule{60}{159.90000000000006}{188}{line}
\Box{62}{160.45000000000007}{73}{\sffamily\fontsize{8.7}{11.5}\selectfont\color{ink}数据供给方多样性}
\Box{140}{160.45000000000007}{14}{\sffamily\bfseries\fontsize{8.6}{11.5}\selectfont\color{teal}诊断}
\Box{158}{160.45000000000007}{27}{\sffamily\fontsize{8.5}{11.5}\selectfont\color{muted}9、10}
\Rule{60}{164.55000000000007}{188}{line}
\Rect{22}{164.80000000000007}{166}{6.5}{gold!9}
\Box{25}{165.80000000000007}{160}{\sffamily\bfseries\fontsize{9.2}{12.5}\selectfont\color{gold}X3 全球连接与网络效应}
\Box{25}{172.10000000000008}{34}{\sffamily\bfseries\fontsize{8.3}{11.3}\selectfont\color{navy}D06 代码影响力}
\Box{62}{172.10000000000008}{73}{\sffamily\fontsize{8.7}{11.5}\selectfont\color{ink}代码仓库影响力}
\Box{140}{172.10000000000008}{14}{\sffamily\bfseries\fontsize{8.6}{11.5}\selectfont\color{gold}北极星}
\Box{158}{172.10000000000008}{27}{\sffamily\fontsize{8.5}{11.5}\selectfont\color{muted}9、17}
\Rule{60}{176.20000000000007}{188}{line}
\Box{62}{176.75000000000009}{73}{\sffamily\fontsize{8.7}{11.5}\selectfont\color{ink}高影响力区域数量}
\Box{140}{176.75000000000009}{14}{\sffamily\bfseries\fontsize{8.6}{11.5}\selectfont\color{gold}诊断}
\Box{158}{176.75000000000009}{27}{\sffamily\fontsize{8.5}{11.5}\selectfont\color{muted}10、11}
\Rule{60}{180.85000000000008}{188}{line}
\Box{62}{181.4000000000001}{73}{\sffamily\fontsize{8.7}{11.5}\selectfont\color{ink}高影响力组织数量}
\Box{140}{181.4000000000001}{14}{\sffamily\bfseries\fontsize{8.6}{11.5}\selectfont\color{gold}诊断}
\Box{158}{181.4000000000001}{27}{\sffamily\fontsize{8.5}{11.5}\selectfont\color{muted}17}
\Rule{60}{185.50000000000009}{188}{line}
\Box{62}{186.0500000000001}{73}{\sffamily\fontsize{8.7}{11.5}\selectfont\color{ink}高影响力项目数量}
\Box{140}{186.0500000000001}{14}{\sffamily\bfseries\fontsize{8.6}{11.5}\selectfont\color{gold}诊断}
\Box{158}{186.0500000000001}{27}{\sffamily\fontsize{8.5}{11.5}\selectfont\color{muted}9}
\Rule{60}{190.1500000000001}{188}{line}
\Box{62}{190.7000000000001}{73}{\sffamily\fontsize{8.7}{11.5}\selectfont\color{ink}高影响力人才数量}
\Box{140}{190.7000000000001}{14}{\sffamily\bfseries\fontsize{8.6}{11.5}\selectfont\color{gold}诊断}
\Box{158}{190.7000000000001}{27}{\sffamily\fontsize{8.5}{11.5}\selectfont\color{muted}8、9}
\Rule{60}{194.8000000000001}{188}{line}
\Box{25}{195.3500000000001}{34}{\sffamily\bfseries\fontsize{8.3}{11.3}\selectfont\color{navy}D07 模型影响力}
\Box{62}{195.3500000000001}{73}{\sffamily\fontsize{8.7}{11.5}\selectfont\color{ink}模型采用度}
\Box{140}{195.3500000000001}{14}{\sffamily\bfseries\fontsize{8.6}{11.5}\selectfont\color{gold}驱动}
\Box{158}{195.3500000000001}{27}{\sffamily\fontsize{8.5}{11.5}\selectfont\color{muted}9}
\Rule{60}{199.4500000000001}{188}{line}
\Box{62}{200.0000000000001}{73}{\sffamily\fontsize{8.7}{11.5}\selectfont\color{ink}模型社区热度}
\Box{140}{200.0000000000001}{14}{\sffamily\bfseries\fontsize{8.6}{11.5}\selectfont\color{gold}驱动}
\Box{158}{200.0000000000001}{27}{\sffamily\fontsize{8.5}{11.5}\selectfont\color{muted}9、17}
\Rule{60}{204.1000000000001}{188}{line}
\Box{62}{204.65000000000012}{73}{\sffamily\fontsize{8.7}{11.5}\selectfont\color{ink}模型生态衍生力}
\Box{140}{204.65000000000012}{14}{\sffamily\bfseries\fontsize{8.6}{11.5}\selectfont\color{gold}北极星}
\Box{158}{204.65000000000012}{27}{\sffamily\fontsize{8.5}{11.5}\selectfont\color{muted}9}
\Rule{60}{208.7500000000001}{188}{line}
\Box{25}{209.30000000000013}{34}{\sffamily\bfseries\fontsize{8.3}{11.3}\selectfont\color{navy}D08 数据集影响力}
\Box{62}{209.30000000000013}{73}{\sffamily\fontsize{8.7}{11.5}\selectfont\color{ink}数据集采用度}
\Box{140}{209.30000000000013}{14}{\sffamily\bfseries\fontsize{8.6}{11.5}\selectfont\color{gold}驱动}
\Box{158}{209.30000000000013}{27}{\sffamily\fontsize{8.5}{11.5}\selectfont\color{muted}9}
\Rule{60}{213.40000000000012}{188}{line}
\Box{62}{213.95000000000013}{73}{\sffamily\fontsize{8.7}{11.5}\selectfont\color{ink}数据社区热度}
\Box{140}{213.95000000000013}{14}{\sffamily\bfseries\fontsize{8.6}{11.5}\selectfont\color{gold}驱动}
\Box{158}{213.95000000000013}{27}{\sffamily\fontsize{8.5}{11.5}\selectfont\color{muted}9、17}
\Rule{60}{218.05000000000013}{188}{line}
\Box{62}{218.60000000000014}{73}{\sffamily\fontsize{8.7}{11.5}\selectfont\color{ink}数据生态复用率}
\Box{140}{218.60000000000014}{14}{\sffamily\bfseries\fontsize{8.6}{11.5}\selectfont\color{gold}北极星}
\Box{158}{218.60000000000014}{27}{\sffamily\fontsize{8.5}{11.5}\selectfont\color{muted}9、12}
\Rule{60}{222.70000000000013}{188}{line}
\Rect{22}{222.95000000000013}{166}{6.5}{violet!9}
\Box{25}{223.95000000000013}{160}{\sffamily\bfseries\fontsize{9.2}{12.5}\selectfont\color{violet}X4 开源善治与公共价值}
\Box{25}{230.25000000000014}{34}{\sffamily\bfseries\fontsize{8.3}{11.3}\selectfont\color{navy}D09 治理成熟度}
\Box{62}{230.25000000000014}{73}{\sffamily\fontsize{8.7}{11.5}\selectfont\color{ink}基金会托管比例}
\Box{140}{230.25000000000014}{14}{\sffamily\bfseries\fontsize{8.6}{11.5}\selectfont\color{violet}北极星}
\Box{158}{230.25000000000014}{27}{\sffamily\fontsize{8.5}{11.5}\selectfont\color{muted}16、17}
\Rule{60}{234.35000000000014}{188}{line}
\Box{62}{234.90000000000015}{73}{\sffamily\fontsize{8.7}{11.5}\selectfont\color{ink}治理主体均衡度}
\Box{140}{234.90000000000015}{14}{\sffamily\bfseries\fontsize{8.6}{11.5}\selectfont\color{violet}诊断}
\Box{158}{234.90000000000015}{27}{\sffamily\fontsize{8.5}{11.5}\selectfont\color{muted}10、16}
\Rule{60}{239.00000000000014}{188}{line}
\Box{62}{239.55000000000015}{73}{\sffamily\fontsize{8.7}{11.5}\selectfont\color{ink}合规许可覆盖率}
\Box{140}{239.55000000000015}{14}{\sffamily\bfseries\fontsize{8.6}{11.5}\selectfont\color{violet}诊断}
\Box{158}{239.55000000000015}{27}{\sffamily\fontsize{8.5}{11.5}\selectfont\color{muted}16}
\Rule{60}{243.65000000000015}{188}{line}
\Box{25}{244.20000000000016}{34}{\sffamily\bfseries\fontsize{8.3}{11.3}\selectfont\color{navy}D10 公共福祉}
\Box{62}{244.20000000000016}{73}{\sffamily\fontsize{8.7}{11.5}\selectfont\color{ink}数字公共品共建度}
\Box{140}{244.20000000000016}{14}{\sffamily\bfseries\fontsize{8.6}{11.5}\selectfont\color{violet}北极星}
\Box{158}{244.20000000000016}{27}{\sffamily\fontsize{8.5}{11.5}\selectfont\color{muted}9、10、17}
\Rule{60}{248.30000000000015}{188}{line}
\Box{62}{248.85000000000016}{73}{\sffamily\fontsize{8.7}{11.5}\selectfont\color{ink}高开放度许可占比}
\Box{140}{248.85000000000016}{14}{\sffamily\bfseries\fontsize{8.6}{11.5}\selectfont\color{violet}诊断}
\Box{158}{248.85000000000016}{27}{\sffamily\fontsize{8.5}{11.5}\selectfont\color{muted}9、16}
\Rule{60}{252.95000000000016}{188}{line}
\Box{62}{253.50000000000017}{73}{\sffamily\fontsize{8.7}{11.5}\selectfont\color{ink}AI 资产透明度}
\Box{140}{253.50000000000017}{14}{\sffamily\bfseries\fontsize{8.6}{11.5}\selectfont\color{violet}诊断}
\Box{158}{253.50000000000017}{27}{\sffamily\fontsize{8.5}{11.5}\selectfont\color{muted}9、16}
\Rule{60}{257.60000000000014}{188}{line}

\Note{22}{266}{166}{表 1-1　用途类型解释见图 1-5。来源：获批第一章候选目录及用户补充指标表。SDG 编号对应图 1-3，仅为目标关联，不是联合国认证、官方指标对应或成效证明；指标定义与数据可得性见附录 A。}
\end{canvas}

\begin{canvas}
\Header{1.3 · 生态对象与关系}{gold}\Footer{09}
\ContinueTitle{gold}{1.3 / 从独立计数到关系观察}{六类对象，构成怎样的开源生态}
\Subhead{59}{gold}{人才与组织共同建设项目，项目连接三类开放资产}
\Body{22}{72}{166}{开源生态既包括参与协作的人与组织，也包括持续开展共同工作的项目，以及被开发、维护和复用的代码、模型与数据。同一个人可参与多个项目，同一项资产也可被多个项目采用。关系能帮助解释单独排名看不到的分工。}
\shade[left color=blue!10,right color=blue!2,rounded corners=2mm,draw=blue!17,line width=.45pt] \Pos{22}{100} rectangle \Pos{102}{124};
\path[fill=blue!9] \Pos{50.8}{124} .. controls \Pos{76.4}{123} and \Pos{82.8}{119} .. \Pos{101}{120} -- \Pos{101}{123} -- \Pos{50.8}{124} -- cycle;
\Box{26}{104}{7}{\sffamily\fontsize{9.5}{14}\selectfont\color{ink}\includegraphics[width=7mm]{.build/assets/D01.pdf}}
\Box{37}{104}{59}{\sffamily\bfseries\fontsize{11}{15}\selectfont\color{blue}人才}
\Box{26}{115}{72}{\sffamily\fontsize{8.6}{12}\selectfont\color{ink}开发者、研究者、维护者与使用者}
\shade[left color=violet!10,right color=violet!2,rounded corners=2mm,draw=violet!17,line width=.45pt] \Pos{108}{100} rectangle \Pos{188}{124};
\path[fill=violet!9] \Pos{136.8}{124} .. controls \Pos{162.4}{123} and \Pos{168.8}{119} .. \Pos{187}{120} -- \Pos{187}{123} -- \Pos{136.8}{124} -- cycle;
\Box{112}{104}{7}{\sffamily\fontsize{9.5}{14}\selectfont\color{ink}\includegraphics[width=7mm]{.build/assets/D02.pdf}}
\Box{123}{104}{59}{\sffamily\bfseries\fontsize{11}{15}\selectfont\color{violet}组织}
\Box{112}{115}{72}{\sffamily\fontsize{8.6}{12}\selectfont\color{ink}企业、高校、基金会与公共机构}
\LinkLine{62}{125}{79}{135}{blue!60}
\Box{23}{128}{46}{\sffamily\fontsize{8}{11}\selectfont\color{blue}贡献、审阅、反馈}
\LinkLine{148}{125}{131}{135}{violet!60}
\Box{143}{128}{45}{\sffamily\fontsize{8}{11}\selectfont\color{violet}投入、治理、支持}
\shade[left color=teal!10,right color=teal!2,rounded corners=2mm,draw=teal!17,line width=.45pt] \Pos{61}{137} rectangle \Pos{149}{159};
\path[fill=teal!9] \Pos{92.68}{159} .. controls \Pos{120.84}{158} and \Pos{127.88}{154} .. \Pos{148}{155} -- \Pos{148}{158} -- \Pos{92.68}{159} -- cycle;
\Box{65}{141}{7}{\sffamily\fontsize{9.5}{14}\selectfont\color{ink}\includegraphics[width=7mm]{.build/assets/D03.pdf}}
\Box{77}{140}{65}{\sffamily\bfseries\fontsize{10}{14}\selectfont\color{teal}项目：围绕共同目标持续协作}
\Box{65}{150}{80}{\sffamily\fontsize{8.5}{12}\selectfont\color{ink}如软件系统、研究工具或公共服务方案}
\draw[muted!35,line width=0.6pt] \Pos{105}{160} .. controls \Pos{105}{166} and \Pos{48}{165} .. \Pos{48}{171};
\shade[left color=blue!10,right color=blue!2,rounded corners=2mm,draw=blue!17,line width=.45pt] \Pos{22}{172} rectangle \Pos{74}{196};
\path[fill=blue!9] \Pos{40.72}{196} .. controls \Pos{57.36}{195} and \Pos{61.52}{191} .. \Pos{73}{192} -- \Pos{73}{195} -- \Pos{40.72}{196} -- cycle;
\Box{25}{176}{6}{\sffamily\fontsize{9.5}{14}\selectfont\color{ink}\includegraphics[width=6mm]{.build/assets/D03.pdf}}
\Box{34}{175}{37}{\sffamily\bfseries\fontsize{10.5}{15}\selectfont\color{blue}代码}
\Box{25}{186}{46}{\sffamily\fontsize{8.1}{11}\selectfont\color{ink}实现功能的程序与工具}
\draw[muted!35,line width=0.6pt] \Pos{105}{160} .. controls \Pos{105}{166} and \Pos{105}{165} .. \Pos{105}{171};
\shade[left color=teal!10,right color=teal!2,rounded corners=2mm,draw=teal!17,line width=.45pt] \Pos{79}{172} rectangle \Pos{131}{196};
\path[fill=teal!9] \Pos{97.72}{196} .. controls \Pos{114.36}{195} and \Pos{118.52000000000001}{191} .. \Pos{130}{192} -- \Pos{130}{195} -- \Pos{97.72}{196} -- cycle;
\Box{82}{176}{6}{\sffamily\fontsize{9.5}{14}\selectfont\color{ink}\includegraphics[width=6mm]{.build/assets/D04.pdf}}
\Box{91}{175}{37}{\sffamily\bfseries\fontsize{10.5}{15}\selectfont\color{teal}模型}
\Box{82}{186}{46}{\sffamily\fontsize{8.1}{11}\selectfont\color{ink}训练后用于预测或生成的模型文件}
\draw[muted!35,line width=0.6pt] \Pos{105}{160} .. controls \Pos{105}{166} and \Pos{162}{165} .. \Pos{162}{171};
\shade[left color=gold!10,right color=gold!2,rounded corners=2mm,draw=gold!17,line width=.45pt] \Pos{136}{172} rectangle \Pos{188}{196};
\path[fill=gold!9] \Pos{154.72}{196} .. controls \Pos{171.36}{195} and \Pos{175.52}{191} .. \Pos{187}{192} -- \Pos{187}{195} -- \Pos{154.72}{196} -- cycle;
\Box{139}{176}{6}{\sffamily\fontsize{9.5}{14}\selectfont\color{ink}\includegraphics[width=6mm]{.build/assets/D05.pdf}}
\Box{148}{175}{37}{\sffamily\bfseries\fontsize{10.5}{15}\selectfont\color{gold}数据}
\Box{139}{186}{46}{\sffamily\fontsize{8.1}{11}\selectfont\color{ink}训练、评测或服务所需资料}
\Box{22}{200}{166}{\sffamily\fontsize{8.3}{12}\selectfont\color{muted}项目组织研发、采用与维护；代码调用模型，数据支持训练或评测，三类资产还可被其他项目复用。}

\Label{22}{210}{166}{用一个例子读懂：共建地方语言公共服务助手（示意）}
\Body{22}{222}{166}{研究者、工程师与本地使用者是\textbf{人才}；大学、企业和基金会是\textbf{组织}。他们围绕服务目标持续协作，形成一个\textbf{项目}。项目可以采用已有\textbf{模型}，共同整理本地语言\textbf{数据}，再开发连接服务的\textbf{代码}。}
\Body{22}{245}{166}{项目不等于单个仓库：一个项目可能维护多个代码仓库、模型版本和数据集。判断其公共价值，还要核验服务成本、维护责任与实际受益。}
\Note{22}{266}{166}{图 1-6　六类对象与关系示意，非实测网络。来源：报告研究框架与\href{\OD}{OpenDigger} 方法材料。人员、组织、项目与资产只有在身份、关系核验后才联读；发布者所在地不自动等于全部贡献者或使用者所在地。}
\end{canvas}

\begin{canvas}
\Header{1.3 · 参与、贡献与影响力}{gold}\Footer{10}
\ContinueTitle{gold}{1.3 / 同一案例，三个观察视角}{参与、贡献与影响力各说明什么}
\Subhead{59}{gold}{从“有人参与”，到“共同完成”，再到“被连接与接续”}
\Body{22}{72}{166}{参与说明谁进入共同工作，贡献关注完成了哪些协作，影响力观察这些工作在关系网络中的位置。AI 生成与自动化改变了活动量的含义，因此要区分参与主体、协作类型与持续性；提交次数不能直接代表能力。}
\Box{22}{102}{166}{\sffamily\bfseries\fontsize{10.4}{15}\selectfont\color{navy}同一案例：项目 P 的兼容性修复，在项目 Q 中继续使用}
\draw[blue!35,line width=1pt,-{Stealth[length=1.4mm]}] \Pos{48}{130} .. controls \Pos{56}{126} and \Pos{65}{134} .. \Pos{75}{130};
\fill[blue!5] \Pos{40}{130} circle[radius=7.7mm];
\shade[top color=white,bottom color=blue!28,shading angle=30,draw=blue!35,line width=.45pt] \Pos{40}{130} circle[radius=6.5mm];
\draw[white,line width=.7pt] \Pos{36.230000000000004}{126.23} arc[start angle=135,end angle=65,radius=5.33mm];
\Box{36.62}{126.62}{6.76}{\sffamily\fontsize{9.5}{14}\selectfont\color{ink}\includegraphics[width=6.76mm]{.build/assets/D01.pdf}}
\Box{21}{140}{38}{\sffamily\bfseries\fontsize{9.2}{13}\selectfont\color{blue}\centering 甲｜发现问题}
\Box{20}{149}{40}{\sffamily\fontsize{8}{11}\selectfont\color{muted}\centering 报告不兼容现象}
\draw[teal!35,line width=1pt,-{Stealth[length=1.4mm]}] \Pos{91}{130} .. controls \Pos{99}{126} and \Pos{108}{134} .. \Pos{118}{130};
\fill[teal!5] \Pos{83}{130} circle[radius=7.7mm];
\shade[top color=white,bottom color=teal!28,shading angle=30,draw=teal!35,line width=.45pt] \Pos{83}{130} circle[radius=6.5mm];
\draw[white,line width=.7pt] \Pos{79.23}{126.23} arc[start angle=135,end angle=65,radius=5.33mm];
\Box{79.62}{126.62}{6.76}{\sffamily\fontsize{9.5}{14}\selectfont\color{ink}\includegraphics[width=6.76mm]{.build/assets/D03.pdf}}
\Box{64}{140}{38}{\sffamily\bfseries\fontsize{9.2}{13}\selectfont\color{teal}\centering 乙｜提交修复}
\Box{63}{149}{40}{\sffamily\fontsize{8}{11}\selectfont\color{muted}\centering 贡献代码改动}
\draw[gold!35,line width=1pt,-{Stealth[length=1.4mm]}] \Pos{134}{130} .. controls \Pos{142}{126} and \Pos{151}{134} .. \Pos{161}{130};
\fill[gold!5] \Pos{126}{130} circle[radius=7.7mm];
\shade[top color=white,bottom color=gold!28,shading angle=30,draw=gold!35,line width=.45pt] \Pos{126}{130} circle[radius=6.5mm];
\draw[white,line width=.7pt] \Pos{122.23}{126.23} arc[start angle=135,end angle=65,radius=5.33mm];
\Box{122.62}{126.62}{6.76}{\sffamily\fontsize{9.5}{14}\selectfont\color{ink}\includegraphics[width=6.76mm]{.build/assets/D09.pdf}}
\Box{107}{140}{38}{\sffamily\bfseries\fontsize{9.2}{13}\selectfont\color{gold}\centering 丙｜审阅验证}
\Box{106}{149}{40}{\sffamily\fontsize{8}{11}\selectfont\color{muted}\centering 测试并完善修改}
\fill[violet!5] \Pos{169}{130} circle[radius=7.7mm];
\shade[top color=white,bottom color=violet!28,shading angle=30,draw=violet!35,line width=.45pt] \Pos{169}{130} circle[radius=6.5mm];
\draw[white,line width=.7pt] \Pos{165.23}{126.23} arc[start angle=135,end angle=65,radius=5.33mm];
\Box{165.62}{126.62}{6.76}{\sffamily\fontsize{9.5}{14}\selectfont\color{ink}\includegraphics[width=6.76mm]{.build/assets/D08.pdf}}
\Box{150}{140}{38}{\sffamily\bfseries\fontsize{9.2}{13}\selectfont\color{violet}\centering 丁｜接续使用}
\Box{149}{149}{40}{\sffamily\fontsize{8}{11}\selectfont\color{muted}\centering 在项目 Q 中复用}
\shade[left color=blue!10,right color=blue!2,rounded corners=2mm,draw=blue!17,line width=.45pt] \Pos{22}{172} rectangle \Pos{74}{215};
\path[fill=blue!9] \Pos{40.72}{215} .. controls \Pos{57.36}{214} and \Pos{61.52}{210} .. \Pos{73}{211} -- \Pos{73}{214} -- \Pos{40.72}{215} -- cycle;
\fill[blue!5] \Pos{30}{181} circle[radius=5.2mm];
\shade[top color=white,bottom color=blue!28,shading angle=30,draw=blue!35,line width=.45pt] \Pos{30}{181} circle[radius=4mm];
\draw[white,line width=.7pt] \Pos{27.68}{178.68} arc[start angle=135,end angle=65,radius=3.28mm];
\Box{27.92}{178.92}{4.16}{\sffamily\fontsize{9.5}{14}\selectfont\color{ink}\includegraphics[width=4.16mm]{.build/assets/D01.pdf}}
\Box{38}{176}{32}{\sffamily\bfseries\fontsize{11}{15}\selectfont\color{blue}参与}
\Box{26}{190}{44}{\sffamily\fontsize{8.8}{12}\selectfont\color{navy}谁持续进入协作？}
\Box{26}{202}{44}{\sffamily\fontsize{8.3}{12}\selectfont\color{muted}人数、留存与角色}
\shade[left color=teal!10,right color=teal!2,rounded corners=2mm,draw=teal!17,line width=.45pt] \Pos{79}{172} rectangle \Pos{131}{215};
\path[fill=teal!9] \Pos{97.72}{215} .. controls \Pos{114.36}{214} and \Pos{118.52000000000001}{210} .. \Pos{130}{211} -- \Pos{130}{214} -- \Pos{97.72}{215} -- cycle;
\fill[teal!5] \Pos{87}{181} circle[radius=5.2mm];
\shade[top color=white,bottom color=teal!28,shading angle=30,draw=teal!35,line width=.45pt] \Pos{87}{181} circle[radius=4mm];
\draw[white,line width=.7pt] \Pos{84.68}{178.68} arc[start angle=135,end angle=65,radius=3.28mm];
\Box{84.92}{178.92}{4.16}{\sffamily\fontsize{9.5}{14}\selectfont\color{ink}\includegraphics[width=4.16mm]{.build/assets/D03.pdf}}
\Box{95}{176}{32}{\sffamily\bfseries\fontsize{11}{15}\selectfont\color{teal}贡献}
\Box{83}{190}{44}{\sffamily\fontsize{8.8}{12}\selectfont\color{navy}完成了哪些共同工作？}
\Box{83}{202}{44}{\sffamily\fontsize{8.3}{12}\selectfont\color{muted}开发、讨论与审阅}
\shade[left color=gold!10,right color=gold!2,rounded corners=2mm,draw=gold!17,line width=.45pt] \Pos{136}{172} rectangle \Pos{188}{215};
\path[fill=gold!9] \Pos{154.72}{215} .. controls \Pos{171.36}{214} and \Pos{175.52}{210} .. \Pos{187}{211} -- \Pos{187}{214} -- \Pos{154.72}{215} -- cycle;
\fill[gold!5] \Pos{144}{181} circle[radius=5.2mm];
\shade[top color=white,bottom color=gold!28,shading angle=30,draw=gold!35,line width=.45pt] \Pos{144}{181} circle[radius=4mm];
\draw[white,line width=.7pt] \Pos{141.68}{178.68} arc[start angle=135,end angle=65,radius=3.28mm];
\Box{141.92}{178.92}{4.16}{\sffamily\fontsize{9.5}{14}\selectfont\color{ink}\includegraphics[width=4.16mm]{.build/assets/D06.pdf}}
\Box{152}{176}{32}{\sffamily\bfseries\fontsize{11}{15}\selectfont\color{gold}影响力}
\Box{140}{190}{44}{\sffamily\fontsize{8.8}{12}\selectfont\color{navy}工作怎样被连接与接续？}
\Box{140}{202}{44}{\sffamily\fontsize{8.3}{12}\selectfont\color{muted}协作关系与结构位置}
\Box{22}{219}{166}{\sffamily\fontsize{8.1}{12}\selectfont\color{muted}图 1-7　协作示意，非实测样本。OpenRank 综合协作关系，观察贡献在网络中的结构性影响。}

\Rule{22}{230}{188}{line}
\Label{22}{236}{166}{主统计期：2025-07-01 至 2026-06-30}
\Note{22}{247}{166}{历史比较使用可比窗口；外部平台年度窗与期后观察分别标注。跨平台身份、自动化账号、地区未知与项目边界按实际核验状态处理，未知不按零计算。}
\Note{22}{262}{166}{来源：获批研究框架与\href{\OD}{OpenDigger／OpenRank 方法材料}。OpenRank 依据协作关系观察网络中的结构性影响；下载、关注、派生与实际公共结果需要分别解释。}
\end{canvas}

\begin{canvas}
\Header{1.3 · 全书证据与阅读路径}{gold}\Footer{11}
\ContinueTitle{gold}{1.3 / 从中国出发的全球视野}{四类证据怎样共同支持决策}
\Subhead{59}{gold}{按问题选择观察视角，再把不同层次的证据联系起来}
\Body{22}{72}{166}{四个研究视角分别组织地区格局、主体影响、技术结构与模型生态。它们共享人才、组织、项目与资产等对象，却回答不同问题。中国的模型、产业与城市实践放在全球接续中观察，全球南方的参与也不能只看资源数量。}
\Box{22}{101}{30}{\sffamily\bfseries\fontsize{10}{14}\selectfont\color{blue}\centering 参与}
\LinkLine{52}{104}{55}{104}{muted!40}
\Box{56}{101}{30}{\sffamily\bfseries\fontsize{10}{14}\selectfont\color{teal}\centering 供给}
\LinkLine{86}{104}{89}{104}{muted!40}
\Box{90}{101}{30}{\sffamily\bfseries\fontsize{10}{14}\selectfont\color{gold}\centering 复用}
\LinkLine{120}{104}{123}{104}{muted!40}
\Box{124}{101}{30}{\sffamily\bfseries\fontsize{10}{14}\selectfont\color{violet}\centering 治理}
\LinkLine{154}{104}{157}{104}{muted!40}
\Box{158}{101}{30}{\sffamily\bfseries\fontsize{10}{14}\selectfont\color{teal}\centering 公共价值}
\Box{22}{110}{166}{\sffamily\fontsize{8.3}{12}\selectfont\color{muted}五段主线贯穿全书；四个视角各有职责，关系核验后才联合解释。}
\Rect{22}{120}{166}{8}{navy}
\Box{25}{121}{31}{\sffamily\bfseries\fontsize{8.8}{12}\selectfont\color{white}章节与视角}
\Box{61}{121}{64}{\sffamily\bfseries\fontsize{8.8}{12}\selectfont\color{white}主要对象与证据}
\Box{129}{121}{56}{\sffamily\bfseries\fontsize{8.8}{12}\selectfont\color{white}可以支持的判断}
\Rect{22}{130}{3}{25}{blue}
\Box{28}{132}{31}{\sffamily\bfseries\fontsize{9.2}{13}\selectfont\color{blue}02 生态地图}
\Box{28}{141}{31}{\sffamily\fontsize{8}{11}\selectfont\color{muted}GOSO}
\Box{61}{132}{64}{\sffamily\bfseries\fontsize{9.3}{13}\selectfont\color{navy}地区、人才与资产}
\Box{61}{142}{64}{\sffamily\fontsize{8.4}{12}\selectfont\color{muted}空间分布、增长、贡献强度、未知覆盖}
\Box{129}{132}{56}{\sffamily\fontsize{8.9}{13.2}\selectfont\color{ink}哪里在增长，增长是否形成持续贡献；识别地区和人才结构。}
\Rule{22}{155}{188}{line}
\Rect{22}{157}{3}{25}{gold}
\Box{28}{159}{31}{\sffamily\bfseries\fontsize{9.2}{13}\selectfont\color{gold}03 影响力排行}
\Box{28}{168}{31}{\sffamily\fontsize{8}{11}\selectfont\color{muted}GOSR}
\Box{61}{159}{64}{\sffamily\bfseries\fontsize{9.3}{13}\selectfont\color{navy}地区、企业、基金会与项目}
\Box{61}{169}{64}{\sffamily\fontsize{8.4}{12}\selectfont\color{muted}协作网络、影响位置、排名与时序}
\Box{129}{159}{56}{\sffamily\fontsize{8.9}{13.2}\selectfont\color{ink}谁承担关键连接与长期维护；识别合作伙伴与结构变化。}
\Rule{22}{182}{188}{line}
\Rect{22}{184}{3}{25}{teal}
\Box{28}{186}{31}{\sffamily\bfseries\fontsize{9.2}{13}\selectfont\color{teal}04 技术全景}
\Box{28}{195}{31}{\sffamily\fontsize{8}{11}\selectfont\color{muted}GOSL}
\Box{61}{186}{64}{\sffamily\bfseries\fontsize{9.3}{13}\selectfont\color{navy}技术类别、项目与工具链}
\Box{61}{196}{64}{\sffamily\fontsize{8.4}{12}\selectfont\color{muted}技术分层、项目联系、维护与演进}
\Box{129}{186}{56}{\sffamily\fontsize{8.9}{13.2}\selectfont\color{ink}模型之外还需要哪些基础工具；识别技术支撑与协作缺口。}
\Rule{22}{209}{188}{line}
\Rect{22}{211}{3}{25}{violet}
\Box{28}{213}{31}{\sffamily\bfseries\fontsize{9.2}{13}\selectfont\color{violet}05 模型生态}
\Box{28}{222}{31}{\sffamily\fontsize{8}{11}\selectfont\color{muted}GOSM}
\Box{61}{213}{64}{\sffamily\bfseries\fontsize{9.3}{13}\selectfont\color{navy}发布者、模型、数据与复用者}
\Box{61}{223}{64}{\sffamily\fontsize{8.4}{12}\selectfont\color{muted}开放供给、衍生网络、许可和使用条件}
\Box{129}{213}{56}{\sffamily\fontsize{8.9}{13.2}\selectfont\color{ink}中国模型怎样进入全球再创造；何种条件支持本地化与普惠。}
\Rule{22}{236}{188}{line}

\Label{22}{244}{166}{第六章：由经验证的发现，形成可检验的行动主张}
\Body{22}{255}{166}{归纳十二项关键洞察，形成分层时间的十大趋势判断与十大战略建议。每项建议交代执行主体、所需条件和观察信号；证据不足时收窄判断。}
\Note{22}{272}{166}{图 1-8　全书证据组织与决策阅读路径。附录 A 同步维护指标定义、数据窗口、映射状态与复核方法。}
\end{canvas}
\end{document}
